\section{Selected tests and their supports}\label{sec:tests}

The next goal is to construct a Banach space that lifts the correction
\(h(s)=(S_s-P)s\) and contains no linear \(c_0\). We use the coordinate
functions \(f_\gamma\) of \eqref{eq:correction-coordinate}, with exactly the finite-dependence and
support assertions \eqref{eq:correction-bounds}--\eqref{eq:correction-support}. The operations below enlarge the family
of tests while retaining quantitative control of where each test can
be nonzero.

\subsection{Two local operations}

Fix a label \(\gamma=(\lambda,j,\xi)\) and its nonempty finite set
\(I=I_\gamma\). A function on \(B\) depending only on \(I\) is
identified with its function on the cube \([-1,1]^I\), by evaluating
at the zero extension of a cube point. Write \(s|_I\) for coordinate
restriction. For a dyadic scale \(\varepsilon=2^{-a}\),
\(a\in\N_0=\{0,1,\ldots\}\), let
\[
 \mathcal Q_{I,\varepsilon}
   =\bigl([-1,1]\cap(\varepsilon/8)\mathbb Z\bigr)^I .
\]
This is a finite grid including both endpoints in each coordinate.
Every point of the cube is within \(\varepsilon/16\), hence within
\(\varepsilon/8\), of a grid point in the supremum norm.

Define the \(2\)-Lipschitz function
\[
 \chi(t)=
 \begin{cases}
 0,&0\leq t\leq1/2,\\
 2t-1,&1/2<t<1,\\
 1,&t\geq1.
 \end{cases}
\]
For a Lipschitz function \(u\) depending on \(I\), and
\(q\in\mathcal Q_{I,\varepsilon}\), put
\begin{align}
 (G_{q,\varepsilon}u)(s)
 &=u(q)+\chi\!\left(\frac{\|s|_I-q\|_\infty}{\varepsilon}\right)
                    (u(s)-u(q)), \label{eq:G}\\
 R_{q,\varepsilon}u&=G_{q,\varepsilon}u-u. \label{eq:R}
\end{align}
Here \(u(q)\) means evaluation at the zero extension of \(q\).
Both operations are linear in \(u\), retain dependence on \(I\), and
\(G_{q,\varepsilon}u\) is the constant \(u(q)\) on the open cylinder
\(\|s|_I-q\|_\infty<\varepsilon/2\).
Thus \(G\) makes \(u\) locally constant and changes it only inside
the radius-\(\varepsilon\) cylinder; \(R\) records that change.
The identity \(G=\mathrm{Id}+R\) will separate fixed parts of a test
from changes occurring at small scales.

\begin{lemma}[Bounds and geometric support]\label{lem:operations}
For every such \(u,q,\varepsilon\),
\begin{align}
 \|R_{q,\varepsilon}u\|_\infty
   &\leq\varepsilon\Lip(u),&
 \Lip(R_{q,\varepsilon}u)&\leq3\Lip(u),&
 \Lip(G_{q,\varepsilon}u)&\leq4\Lip(u), \label{eq:operation-lip}\\
 \|G_{q,\varepsilon}u\|_\infty
   &\leq\|u\|_\infty,&
 \|R_{q,\varepsilon}u\|_\infty&\leq2\|u\|_\infty. \label{eq:operation-sup}
\end{align}
Let
\[
 K_I(u)=\overline{\{s|_I:s\in B,\ u(s)\ne0\}}^{\,[-1,1]^I}.
\]
In the next display write \(G=G_{q,\varepsilon}\) and
\(R=R_{q,\varepsilon}\).
With \(K_\varepsilon=\{p:\operatorname{dist}_\infty(p,K)\leq
\varepsilon\}\) and \(\varnothing_\varepsilon=\varnothing\), one has
\begin{equation}
 K_I(Gu),K_I(Ru)\subset K_I(u)_\varepsilon,\qquad
 K_I(Ru)\subset\overline B_\infty(q,\varepsilon).             \label{eq:operation-support}
\end{equation}
\end{lemma}

\begin{proof}
Set \(a(s)=1-\chi(\|s|_I-q\|_\infty/\varepsilon)\). Then
\[
 R_{q,\varepsilon}u(s)=a(s)(u(q)-u(s)),\qquad
 0\leq a\leq1,\qquad \Lip(a)\leq2/\varepsilon.
\]
If \(a(s)\ne0\), then \(\|s|_I-q\|_\infty<\varepsilon\);
the cylindrical Lipschitz bound gives
\(|u(q)-u(s)|\leq\varepsilon\Lip(u)\). This proves the first
supremum estimate in \eqref{eq:operation-lip}.

To compare \(Ru(s)\) and \(Ru(t)\), there is nothing to prove if both
values of \(a\) are zero. Otherwise interchange \(s,t\), if needed,
so that \(a(t)\ne0\). The identity
\[
 Ru(s)-Ru(t)
 =a(s)(u(t)-u(s))+(a(s)-a(t))(u(q)-u(t))
\]
bounds the two terms by \(\Lip(u)\|s-t\|_D\) and
\((2/\varepsilon)\|s-t\|_D\,\varepsilon\Lip(u)\), respectively.
Thus \(\Lip(Ru)\leq3\Lip(u)\), and \(Gu=u+Ru\) gives
\(\Lip(Gu)\leq4\Lip(u)\). The formula for \(G\) is a convex
combination of \(u(q)\) and \(u(s)\), so it does not increase the
supremum norm. The bound for \(R=G-u\) follows.

For the support assertions, if a new output is nonzero at \(s\) while
\(u(s)=0\), then \(u(q)\ne0\) and
\(\|s|_I-q\|_\infty<\varepsilon\). If \(u(s)\ne0\), its projection
already lies in \(K_I(u)\). This covers both \(G\) and \(R\).
Taking closures gives their inclusion in \(K_I(u)_\varepsilon\);
the latter set is closed because the cube is compact. In addition,
\(Ru(s)=0\) whenever \(\|s|_I-q\|_\infty\geq\varepsilon\),
which proves its ball inclusion. If \(u=0\), all statements are
immediate.
\end{proof}

\subsection{Scale budgets and tail factors}

For \(h\geq1\), let \(\mathcal T_h\) consist of functions obtained
from a single \(f_\gamma\), \(\gamma=(\lambda,j,\xi)\), by \(m\leq h\)
successive operations \(G\) or \(R\) on \(I_\gamma\). A representation
with \(m=0\) is allowed. For \(m\geq1\), the scales in their order of
application must satisfy
\begin{equation}
 \varepsilon_1\geq\varepsilon_2\geq\cdots\geq\varepsilon_m>0,
 \qquad
 \sum_{k=1}^m\varepsilon_k<\frac{\ell_j}{2}.                 \label{eq:scale-rule}
\end{equation}
The sum is zero for \(m=0\). Each scale is dyadic and each center is
in its corresponding finite grid. We refer to \(\gamma\) as a
representing label, since a function may have more than one
representation. The strict inequality leaves a positive amount of
budget for a further sufficiently small operation. The order will
ensure that every operation after a small one has at most that scale.

For a continuous function on \(B\), let
\(\supp_B u=\overline{\{s\in B:u(s)\ne0\}}^{\,B}\).
All support statements refer to these concrete generated functions.
Later, molecule pairings will ignore additive constants, but we do not
replace a generated function by another representative modulo constants:
that could change its support and its constant value on a cylinder.

\begin{lemma}[Propagation of the correction support]\label{lem:support}
Suppose \(u\in\mathcal T_h\) has a representation based at
\(\gamma=(\lambda,j,\xi)\) with total scale
\(e=\sum_{k=1}^m\varepsilon_k\). On \(\supp_B u\),
\begin{equation}
 \begin{cases}
 \|\pi_j s-\xi\|_\infty\leq2d_j+e,&j\geq1,\\
 r_\lambda(s)\geq\ell_j-e,&j<J_\lambda.
 \end{cases}                                                  \label{eq:support-sharp}
\end{equation}
In particular, on the same support,
\begin{equation}
 \begin{cases}
 \|\pi_j s-\xi\|_\infty\leq2d_j+\ell_j/2,&j\geq1,\\
 r_\lambda(s)\geq\ell_j/2,&j<J_\lambda.
 \end{cases}                                                  \label{eq:support-coarse}
\end{equation}
\end{lemma}

\begin{proof}
Start with \eqref{eq:correction-support}, which holds on the relative support of \(f_\gamma\)
as well as its nonzero set, because its two inequalities describe
closed subsets of \(B\). Suppose the claim has been proved for the
old function and add an operation of scale \(\varepsilon\).
If the new function is nonzero at \(s\) and the old one is zero there,
the proof of Lemma~\ref{lem:operations} gives an old nonzero value
at the zero extension of \(q\), with
\(\|s|_{I_\gamma}-q\|_\infty<\varepsilon\). Otherwise the old
support already contains \(s\). Both \(\pi_j\) and \(r_\lambda\)
depend only on \(I_\gamma\); coordinate restriction has norm one,
and \(r_\lambda\) is \(1\)-Lipschitz on that cube. The first upper
bound therefore increases by at most \(\varepsilon\), and the
second lower bound decreases by at most \(\varepsilon\).
Induction proves \eqref{eq:support-sharp} on the nonzero set.
Taking relative closure preserves the closed inequalities. Finally,
\eqref{eq:scale-rule} gives \eqref{eq:support-coarse}.
\end{proof}

For \(N\geq1\), define the coordinate tail
\begin{equation}
 \tau_N(s)=\sup_{k>N}|s_k|,\qquad
 b_{c,N}(s)=\min\{1,(2-\tau_N(s)/c)_+\}\quad(c>0),
                                                               \label{eq:tail}
\end{equation}
where \(a_+=\max\{a,0\}\). The function \(\tau_N\) is
\(1\)-Lipschitz, and \(\tau_N(s)\to0\) for each \(s\in B\).
Both \(b_{c,N}\) and \(b_{c,N}-1\) have supremum norm at most one
and Lipschitz constant at most \(1/c\). Their relative supports satisfy
\begin{equation}
 \supp_B b_{c,N}\subset\{\tau_N\leq2c\},\qquad
 \supp_B(b_{c,N}-1)\subset\{\tau_N\geq c\}.                  \label{eq:tail-support}
\end{equation}
For example, the first function vanishes when \(\tau_N\geq2c\),
and the second vanishes when \(\tau_N\leq c\); closure gives the
displayed weak inequalities.

Let \(\mathcal V_h\) consist of all products
\begin{equation}
                         u\prod_{r=1}^l\beta_r              \label{eq:Vclass}
\end{equation}
where \(u\in\mathcal T_h\) has a representing label with \(j\leq h\),
\(0\leq l\leq h\), and each \(\beta_r\) is \(b_{c_r,N_r}\) or
\(b_{c_r,N_r}-1\), with \(N_r\geq1\) and
\[
                       c_r\in\{2^{-a}:0\leq a\leq h\}.
\]
The empty product is one. The support of a product is contained in
the support of its cylindrical factor \(u\), so every such test
inherits \eqref{eq:support-coarse}. We will use the countable list of
classes
\begin{equation}
               \mathcal T_1,\mathcal V_1,
               \mathcal T_2,\mathcal V_2,\ldots.             \label{eq:classes}
\end{equation}

\begin{lemma}[Bounds in each fixed class]\label{lem:class-bounds}
With \(A,L_0\) from \eqref{eq:base-bounds}, the following bounds are uniform over all
tests of the indicated class:
\begin{equation}
 \begin{array}{c|cc}
 &\|u\|_\infty&\Lip(u)\\ \hline
 u\in\mathcal T_h&2^hA&4^hL_0\\
 u\in\mathcal V_h&2^hA&4^hL_0+h4^hA.
 \end{array}                                                  \label{eq:class-bounds}
\end{equation}
\end{lemma}

\begin{proof}
Iterating \eqref{eq:operation-lip}--\eqref{eq:operation-sup}
gives the first row. For the second, every tail factor has supremum
norm at most one and Lipschitz constant at most \(2^h\). A product
of at most \(h\) such factors has Lipschitz constant at most
\(h2^h\), by telescoping the product difference. The product rule
then gives
\[
 \Lip\!\left(u\prod_r\beta_r\right)
 \leq4^hL_0+(2^hA)(h2^h)=4^hL_0+h4^hA.
\]
Its supremum norm is at most \(2^hA\).
\end{proof}

We have two useful consequences of the support bounds. For a fixed
bounded set of bands \(j\), the second line of
\eqref{eq:support-coarse} gives a positive radial threshold whenever
\(j<J_\lambda\). For \(j\to\infty\), the first line confines the
first \(j\) coordinates to a radius
\(2d_j+\ell_j/2\to0\). Section~\ref{sec:no-c0} will use precisely
these alternatives, after the next section extends the tests from
finite molecules to the completed Banach space.
